\section*{الفصل \ref{ch:b3:genomics} --- الجينوميات والتسلسل}

\begin{solution}{exo:b3:genomics:1}
النوكليوتيد \omterm{def:b3:genomics:sanger}{ثنائي منزوع الأكسجين} لا هيدروكسيل 3$'$ له، فلا تُصنع
رابطة فوسفودايستر إلى النوكليوتيد التالي وتتوقف السلسلة. ولو لم
توجد إلا الصور ثنائية منزوعة الأكسجين لتوقفت كل سلسلة عند الموضع
الأول؛ ولو لم توجد إلا الصور السوية لما توقفت أي سلسلة. فالمزيج
يجعل الإنهاء حدثًا عشوائيًا عند كل موضع، فتكوّن النواتج سلّمًا كاملًا،
طولًا واحدًا لكل قاعدة من القالب.
\end{solution}

\begin{solution}{exo:b3:genomics:2}
$4\times 10^{8}\times \qty{300}{bp} = 1.2\times 10^{11}$ قاعدة. وللإنسان:
$1.2\times 10^{11}/3.2\times 10^{9} \approx 38\times$. وللبكتيريا:
$1.2\times 10^{11}/5\times 10^{6} = \num{24000}\times$.
\end{solution}

\begin{solution}{exo:b3:genomics:3}
الكونتيغ متتالية متصلة مجمَّعة من قراءات متداخلة؛ \omterm{def:b3:genomics:assembly}{والسقالة} مجموعة
مرتبة موجَّهة من الكونتيغات بينها فجوات مقدَّرة
الحجم؛ ومقياس \omterm{def:b3:genomics:assembly}{N50} هو طول الكونتيغ الذي تبلغ عنده الكونتيغات المرتبة
نصف القواعد المجمَّعة. وهنا تحوي الكونتيغات العشرة بطول \qty{8}{Mb}
وحدها \qty{80}{Mb}، وهو يتجاوز نقطة النصف \qty{50}{Mb}
(فالسابع يبلغ \qty{56}{Mb}): فمقياس \omterm{def:b3:genomics:assembly}{N50} $= \qty{8}{Mb}$.
\end{solution}

\begin{solution}{exo:b3:genomics:4}
\omterm{prop:b3:genomics:sizes}{حجم الجينوم} لا يتبع تعقيد الكائن: فللبصل خمسة
أضعاف DNA الإنسان، ولسمكة الرئة أربعون ضعفًا؛ وتختلف الخميرة والدودة
عشرة أضعاف في DNA مع تقارب عدد المورّثات. والفائض
غير مشفِّر: عناصر قافزة وبقاياها، وإنترونات،
وتكرارات ساتلة.
\end{solution}

\begin{solution}{exo:b3:genomics:5}
$e^{-c} = 10^{-6}$ يعطي $c = 6\ln 10 = 13.8$. وعند $c = 8$: $N =
cG/L = 8\times 5\times 10^{6}/150 \approx \num{267000}$ \omterm{def:b3:genomics:ngs}{قراءة}، و$\theta
= 30/150 = 0.2$، والكونتيغات $= N e^{-6.4} = \num{267000}\times 0.00166
\approx 440$.
\end{solution}

\begin{solution}{exo:b3:genomics:6}
$P(X < 8) = P(X \le 7) \approx P\bigl(Z < (7.5 - 15)/2.74\bigr) =
P(Z < -2.74) \approx 0.003$. وعند $30\times$ يُرى كلا أليلي
متغاير اللواقح مرات كثيرة دائمًا تقريبًا، والتغطية غير متساوية
(فالمناطق الغنية بالقاعدتين G وC تنال قراءات أقل، فترى بعض المواضع نصف المتوسط)،
ويبقى من القراءات ما يكفي لفصل أليل حقيقي عن أخطاء
معدلها $10^{-3}$.
\end{solution}

\begin{solution}{exo:b3:genomics:7}
كل \omterm{def:b3:genomics:ngs}{قراءة} بطول \qty{150}{bp} من داخل التكرار متطابقة أيًا كانت
النسخة التي جاءت منها، فيبني المجمِّع عقدة واحدة بطول \qty{6}{kb}
يدخلها $500$ مسلك ويخرج منها $500$؛ ولا يستطيع أن يعرف أي مدخل
يقترن بأي مخرج، فينكسر \omterm{def:b3:genomics:assembly}{التجميع} إلى $500$ فجوة، ويحضر
التكرار مرة واحدة بتغطية تبلغ $500$ ضعف المتوسط. والقراءات الأطول
من التكرار مع جناحين فريدين على جانبيه --- أي \qty{8}{kb} أو
أكثر --- تعبر كل نسخة فتحلّها.
\end{solution}

\begin{solution}{exo:b3:genomics:8}
$4.8\times 10^{7}/1000 = \num{48000}$ اختلاف مشفِّر؛ وثلثها،
أي نحو \num{16000}، يغيّر البروتين.
\end{solution}

\begin{solution}{exo:b3:genomics:9}
الإيجابيات الكاذبة المتوقعة $10^{6}\times 5\times 10^{-8} = 0.05$. وخطر
نسبي قدره $1.15$ في مرض تواتره \qty{2}{\%} يغيّر بضع حالات
في الألف؛ وألف شخص يضمون نحو عشرين حالة، وهي أقل بكثير
من أن تُرى (فالقدرة تنمو بعدد الحالات وبمربع
الأثر). أما للفرد فإن $0.3$ نقطة مئوية
لا تعني شيئًا. ولكن الأليل يعلّم مورّثة يغيّر تغيرٌ متواضع في
نشاطها خطرَ المرض --- وقد تكون المورّثة ومسلكها هدفًا
دوائيًا لتثبيطه التام أثر كبير.
\end{solution}

\begin{solution}{exo:b3:genomics:10}
جمهرية وراثية: ففي البكتيريا، وجمهراتها هائلة، يزيل الانتقاء
حتى الإقحامات الطفيفة الكلفة؛ وفي الثدييات، وجمهراتها الفعالة
صغيرة، يدع الانجراف DNA غير المشفِّر الطفيف الضرر
يتراكم --- ويدعم ذلك الارتباط العكسي عبر السلالات
بين \omterm{prop:b3:genomics:sizes}{حجم الجينوم} وحجم الجمهرة، وتضعفه استثناءات.
وبنيوية: فجينومات حقيقيات النوى تغزوها عناصر قافزة تطهّرها
البكتيريا، لانحيازها إلى الحذف ولانعدام ملاذ منصّف
للعناصر الأنانية --- ويدعم ذلك ارتباط \omterm{prop:b3:genomics:sizes}{حجم الجينوم}
بمحتواه من العناصر القافزة. وتنظيمية: فالتطور المعقد يحتاج إلى
DNA منظِّم أكثر --- وتدعم ذلك العناصر غير المشفِّرة المحفوظة
حول المورّثات التطورية، ويضعفه أن
\qty{8}{\%} فقط من الجينوم البشري يُظهر أي انتقاء أصلًا، فمعظم
DNA غير المشفِّر ليس تنظيميًا.
\end{solution}

\begin{solution}{exo:b3:genomics:11}
تقع البدايات من النوعين بكثافة كلية $\rho = (N_{1} +
N_{2})/G$. وتكون \omterm{def:b3:genomics:ngs}{قراءة} طولها $L_{i}$ أقصى كونتيغها يمينًا إذا لم
تبدأ \omterm{def:b3:genomics:ngs}{قراءة} من أي نوع في $L_{i} - T$ قاعدة بعدها:
باحتمال $e^{-\rho(L_{i} - T)}$. فتكون الكونتيغات $= N_{1}
e^{-\rho(L_{1} - T)} + N_{2} e^{-\rho(L_{2} - T)}$. \omterm{def:b3:genomics:ngs}{والقراءة} الطويلة
تكون الأقصى يمينًا باحتمال أصغر أسّيًا من \omterm{def:b3:genomics:ngs}{القراءة} القصيرة،
كما أن كل \omterm{def:b3:genomics:ngs}{قراءة} طويلة تزيل احتمال فجوة على امتداد طولها
كله؛ أما العدد نفسه من القواعد في قراءات قصيرة فيضيف نقاط بدء
كثيرة لكن كل نافذة محمية قصيرة، فتغلب القراءات الطويلة.
\end{solution}

\begin{solution}{exo:b3:genomics:12}
تتنبأ مضاعفتا الجينوم الكامل بأن النمط الرباعي يعم
الجينوم: رباعيات من قطع صبغية متماثلة متوازية
تحمل أسر المورّثات نفسها بالترتيب نفسه، مع تأريخ
المضاعفتين إلى الزمن نفسه في كل الأسر؛ وبعنقود واحد
في البرمائيات الرأسية و\emph{Ciona}، وقد تباعدتا قبل الحدثين؛ وفي
اللمبريات، وقد تباعدت حولهما، بحالة وسطى أو
مشتقة استقلالًا. أما المضاعفتان المستقلتان لعنقود Hox
وحده فتتنبآن بتماثلات متوازية لعنقود Hox وحده، وبجيران لهم
تواريخ مختلفة، وبتواريخ مضاعفة تختلف بين
الأسر. والتماثلات المتوازية العامة للجينوم والتأريخ المشترك، كما
رُصدا، يرجّحان مضاعفة الجينوم الكامل.
\end{solution}

\begin{solution}{pb:b3:genomics:1}
\textbf{1.} $0.002\times 6\times 10^{8} = 1.2\times 10^{6}$ \omterm{def:b3:genomics:ngs}{قراءة}؛
و$c = 1.2\times 10^{6}\times 150/4.6\times 10^{6} \approx 39$.
\textbf{2.} $e^{-39} \approx 10^{-17}$: أي لا قاعدة غير مغطاة عمليًا.
\textbf{3.} $\theta = 0.2$؛ والكونتيغات $= 1.2\times 10^{6} e^{-31} \approx
0$: أي كونتيغ واحد نظريًا، والفجوات عند التكرارات فقط.
\textbf{4.} $c = \num{30000}\times 150/4.6\times 10^{6} = 0.98$؛
وغير المغطى $e^{-0.98} = 0.38$، أي نحو \qty{1.7}{Mb}؛ والكونتيغات $= \num{30000}
\times e^{-0.78} \approx \num{13700}$.
\textbf{5.} الأوبرونات السبعة تعطي قراءات متطابقة فتنهار في كونتيغ
واحد بطول \qty{5}{kb}، تدخله سبعة أجنحة يسرى فريدة وتخرج منه
سبعة أجنحة يمنى لا يُعرف اقترانها: أي $14$ نهاية كونتيغ، وسبع
فجوات.
\textbf{6.} نحو \num{4600} مورّثة؛ والمشفِّر $0.88\times 4.6\times 10^{6}
= 4.0\times 10^{6}$ زوج قواعد، ومتوسط المورّثة \qty{880}{bp}.
\textbf{7.} $0.998\times 6\times 10^{8}\times 150/3.2\times 10^{9}
\approx 28\times$.
\textbf{8.} نحو $14$ \omterm{def:b3:genomics:ngs}{قراءة} لكل أليل.
\textbf{9.} $28\times 10^{-3}/3 \approx 0.009$ \omterm{def:b3:genomics:ngs}{قراءة} تُظهر قاعدة
خاطئة بعينها --- واحتمال ثلاث أو أكثر نحو $10^{-7}$ --- و
\qty{20}{\%} من $28$ هي $5.6$ \omterm{def:b3:genomics:ngs}{قراءة}؛ فالخطأ لا يجتاز أيًا من المحكّين،
بينما يُتوقع الأليل المتغاير الحقيقي في $14$ \omterm{def:b3:genomics:ngs}{قراءة}.
\textbf{10.} $3.2\times 10^{9}/1500 \approx 2.1\times 10^{6}$ موضع
متغاير اللواقح؛ وفي التسلسل المشفِّر $4.8\times 10^{7}/1500 \approx
\num{32000}$.
\textbf{11.} \omterm{def:b3:genomics:ngs}{القراءة} بطول \qty{150}{bp} من عنصر Alu تطابق كثيرًا من
النسخ البالغة $1.1$ مليون مطابقةً متقاربة، فتوضع على النسخة الخاطئة
أو تُعطى ثقة محاذاة منخفضة. ثم تبدو الفروق بين النسخ
متغايرات متغايرة اللواقح، وتختفي المتغايرات الحقيقية بينها:
فتُرشَّح استدعاءات المتغايرات داخل Alu عادةً، وتكون هذه
العناصر بقعًا عمياء في تسلسل القراءات القصيرة.
\textbf{12.} $1.2\times 10^{-8}\times 6.4\times 10^{9} \approx 77$ طفرة
جديدة. وينبغي أن تُظهر نحو $14$ \omterm{def:b3:genomics:ngs}{قراءة} المتغايرَ في الطفل؛ لكن
إذا لم يُغطَّ موضع الوالد إلا بخمس قراءات، فات الأليل متغاير اللواقح
باحتمال $2^{-5} = 3\,\%$ واستُدعي المتغاير الموروث
جديدًا --- فلا بد إذن من تسلسل الوالدين بعمق عمق
الطفل.
\textbf{13.} $4.8\times 10^{7}\times 100/150 = 3.2\times 10^{7}$ \omterm{def:b3:genomics:ngs}{قراءة}،
أي أقل عشرين ضعفًا من الجينوم: فتهبط كلفة التسلسل
عشرين ضعفًا، وهو أكثر من كلفة خطوة الالتقاط.
\textbf{14.} $1.1\times 10^{6}\times 300 = 3.3\times 10^{8}$ زوج قواعد، أي نحو
\qty{10}{\%} من الجينوم؛ و\qty{10}{\%} من القراءات، أي نحو $6\times
10^{7}$.
\textbf{15.} $\num{60000}\times 1500 = 9\times 10^{7}$ زوج قواعد؛ و$9\times
10^{7}/1.6\times 10^{10} = 0.56\,\%$.
\textbf{16.} $0.012\times 2.9\times 10^{9} = 3.5\times 10^{7}$
تبدّل، أي $1.7\times 10^{7}$ لكل سلالة؛ والمعدل $1.7\times
10^{7}/(2.9\times 10^{9}\times 6.5\times 10^{6}) = 9\times 10^{-10}$
لكل قاعدة في السنة، و$2.3\times 10^{-8}$ لكل جيل من 25 سنة ---
أي نحو ضعف معدل الأنساب البالغ $1.2\times 10^{-8}$، وهو تفاوت
معروف يوحي بأزمنة أجيال أطول في الماضي أو بانفصال أقدم.
\textbf{17.} أربع نسخ من \num{16000} تكون \num{64000}؛ وبقيت \num{20000}،
فمن النسخ الزائدة البالغة \num{48000} لم ينجُ إلا \num{4000}:
أي فُقد \qty{92}{\%} من المضاعفات.
\textbf{18.} للدودة والإنسان عدد المورّثات نفسه؛ والتعقيد
في كيفية استعمالها --- فالتشذيب البديل يضاعف المورّثة الواحدة
إلى آلاف البروتينات (\num{38016} في \emph{Dscam})، والتنظيم
يجمع عوامل الاستنساخ والمعزِّزات بطرق خاصة بكل نمط
خلوي، وتضيف RNA غير المشفِّرة طبقات، وتتفاعل البروتينات في شبكات.
\textbf{19.} عناصر منظِّمة (معزِّزات، ومروّجات، وعوازل)،
ومورّثات \omterm{def:b3:rna-regulation:ncrna}{RNA غير مشفِّر}، وإشارات تشذيب وتوطين، ومناشئ
ومتتاليات بنيوية. واختبر عنصرًا معزِّزًا مرشحًا بوضعه قبل
مورّثة كاشفة في جنين نقيل والنظر إلى أين يُعبَّر عن
الكاشفة، ثم بحذف العنصر من الجينوم بمنظومة
CRISPR وقياس المورّثات المجاورة.
\textbf{20.} $0.05/10^{6} = 5\times 10^{-8}$؛ وعند $p < 10^{-5}$، $10^{6}
\times 10^{-5} = 10$ إيجابيات كاذبة متوقعة.
\textbf{21.} أرجحية غير الحامل $0.009/0.991 = 0.00908$؛ وأرجحية نسخة واحدة
$0.00954$، فالخطر $0.945\,\%$؛ وأرجحية نسختين $0.01001$، فالخطر
$0.99\,\%$.
\textbf{22.} تجمع الدرجة، على امتداد الأليلات البالغة $200$، عددَ نسخ
الخطر التي يحملها الشخص موزونةً بلوغاريتم نسبة أرجحية كل أليل.
ومتوسط العدد $200\times 2\times 0.3 = 120$ وانحرافه المعياري
$\sqrt{200\times 2\times 0.3\times 0.7} \approx 9$؛ فأعلى \qty{1}{\%}
يحمل نحو $21$ أليلًا فوق المتوسط، و$1.05^{21} \approx
2.8$: أي قرابة ثلاثة أضعاف أرجحية الشخص المتوسط، من أليلات
يغيّر كل منها الخطر \qty{5}{\%} وحده.
\textbf{23.} الأليلات داخل كتلة من عشرات الكيلوقواعد تُورَّث
معًا (عدم الاتزان الارتباطي)، فيُظهر أي منها الارتباط نفسه؛
والموضع المنمَّط ليس إلا الموضع الموجود على الشريحة. ولإيجاد
المتغاير السببي: سلسِل المنطقة في الحالات والشواهد، وضيّق
المجموعة إحصائيًا، ثم اختبر كل مرشح --- بمقايسات كاشفة
لنشاط كل أليل معزِّزًا، وبتحرير كل أليل في خلايا
وقياس تعبير المورّثات المجاورة، وبالتوضّع المشترك مع إشارات
الصفات الكمية للتعبير.
\textbf{24.} $1 - e^{-0.5} = 39\,\%$ من الجينوم يُقرأ. وهو يعطي
عينة مباشرة من جمهرة في زمن معلوم من الماضي --- تواترات
أليلية قبل الهجرات والاختلاط اللاحقة، وطول قطع
النياندرتال (وهي طويلة، لأن أجيالًا قليلة من إعادة التركيب
كسرتها) ومن ثم تأريخًا للاختلاط --- وهو ما لا تستطيع الجينومات
الحديثة إلا أن تستنتجه بالنماذج.
\textbf{25.} نحو \num{13700} كونتيغ في \omterm{def:b3:genomics:assembly}{التجميع} الأولي؛ و$28\times$ لكل
جينوم أحادي الصيغة، أي نحو $14$ \omterm{def:b3:genomics:ngs}{قراءة} لكل أليل؛ ونحو $77$ طفرة جديدة
في طفل.
\end{solution}
